\section{The detectability lemma and its converse}

Let $h_0,\dots,h_{m-1}$ be orthogonal projections on a complete inner
product space, write $Q_j=1-h_j$, and for a vector $\psi$ put
$\psi_k=Q_{k-1}\cdots Q_0\psi$. The detectability lemma bounds the energy
of $\varphi=\psi_m$ with respect to $H=\sum_jh_j$ by the mass that the
product $Q_{m-1}\cdots Q_0$ removes from $\psi$. It is Lemma~lem:dl of
\cite[lines 1508--1519]{Arad2016RigorousRG}, proved in
\cite{Anshu2016Simple}; its converse is Lemma~lem:dl-converse of
\cite[lines 1528--1537]{Arad2016RigorousRG}.

\begin{definition}[Product of complementary projections]
    \label{def:complement_product}
    \lean{ContinuousLinearMap.complementProd}
    \leanok
    The product $Q_{k-1}\cdots Q_0$, in which $Q_0=1-h_0$ acts first.
\end{definition}

\begin{theorem}[Projection identities]
    \label{thm:star_projection_identities}
    \lean{ContinuousLinearMap.IsStarProjection.apply_apply_self,
        ContinuousLinearMap.IsStarProjection.inner_apply_left,
        ContinuousLinearMap.IsStarProjection.inner_apply_self,
        ContinuousLinearMap.IsStarProjection.re_inner_apply_self,
        ContinuousLinearMap.IsStarProjection.norm_sq_eq,
        ContinuousLinearMap.IsStarProjection.norm_apply_le}
    \leanok
    For an orthogonal projection $p$ and vectors $x,y$:
    $p(px)=px$, $\langle px,y\rangle=\langle x,py\rangle$,
    $\langle x,px\rangle=\|px\|^2$, $\|x\|^2=\|px\|^2+\|x-px\|^2$ and
    $\|px\|\le\|x\|$.
\end{theorem}

\begin{proof}
    \leanok
    The first two hold because $p$ is idempotent and self-adjoint, and they
    give $\langle x,px\rangle=\langle px,px\rangle$ and
    $\langle px,x-px\rangle=0$.
\end{proof}

\begin{theorem}[Removed mass]
    \label{thm:complement_product_removed_mass}
    \lean{ContinuousLinearMap.complementProd_succ_apply,
        ContinuousLinearMap.norm_sq_complementProd_succ,
        ContinuousLinearMap.re_inner_sum_apply_eq,
        ContinuousLinearMap.sum_norm_sq_apply_complementProd}
    \leanok
    \uses{def:complement_product, thm:star_projection_identities}
    For every $k$, $\psi_{k+1}=\psi_k-h_k\psi_k$ and
    $\|\psi_{k+1}\|^2=\|\psi_k\|^2-\|h_k\psi_k\|^2$. Hence
    $\sum_{k<n}\|h_k\psi_k\|^2=\|\psi\|^2-\|\psi_n\|^2$. Moreover
    $\langle y,(\sum_{i\in s}h_i)y\rangle=\sum_{i\in s}\|h_iy\|^2$ for every
    finite set $s$ and vector $y$.
\end{theorem}

\begin{proof}
    \leanok
    Apply Pythagoras to $\psi_k=h_k\psi_k+Q_k\psi_k$ and sum over $k$.
\end{proof}

\begin{theorem}[One term of the energy]
    \label{thm:detectability_one_term}
    \lean{ContinuousLinearMap.norm_apply_complementProd_le}
    \leanok
    \uses{thm:complement_product_removed_mass}
    If $i<n$, then $\|h_i\psi_n\|\le\sum_k\|h_k\psi_k\|$, the sum over the
    indices $i<k<n$ for which $h_k$ does not commute with $h_i$.
\end{theorem}

\begin{proof}
    \leanok
    By induction on $n$. For $n=i+1$, $h_i\psi_{i+1}=h_iQ_i\psi_i=0$. If
    $h_n$ commutes with $h_i$, then $h_i\psi_{n+1}=Q_nh_i\psi_n$, whose norm
    is at most $\|h_i\psi_n\|$. Otherwise
    $h_i\psi_{n+1}=h_i\psi_n-h_ih_n\psi_n$, whose norm is at most
    $\|h_i\psi_n\|+\|h_n\psi_n\|$.
\end{proof}

\begin{theorem}[Detectability lemma]
    \label{thm:detectability_lemma}
    \lean{ContinuousLinearMap.re_inner_sum_complementProd_le}
    \leanok
    \uses{thm:detectability_one_term}
    Suppose each $h_i$ fails to commute with at most $g$ of the other
    $h_j$, $j<m$. Then
    $\langle\varphi|H|\varphi\rangle\le g^2(\|\psi\|^2-\|\varphi\|^2)$.
\end{theorem}

\begin{proof}
    \leanok
    Let $A_i$ be the set of $k$ with $i<k<m$ for which $h_k$ does not
    commute with $h_i$, and $a_k=\|h_k\psi_k\|$. By the previous theorem
    and the Cauchy--Schwarz inequality,
    $\|h_i\varphi\|^2\le|A_i|\sum_{k\in A_i}a_k^2\le g\sum_{k\in A_i}a_k^2$.
    Each $k$ lies in $A_i$ for at most $g$ indices $i$, so
    $\langle\varphi|H|\varphi\rangle=\sum_i\|h_i\varphi\|^2
        \le g^2\sum_{k<m}a_k^2=g^2(\|\psi\|^2-\|\varphi\|^2)$.
\end{proof}

\begin{theorem}[Detectability lemma, normalized form]
    \label{thm:detectability_lemma_normalized}
    \lean{ContinuousLinearMap.norm_sq_complementProd_le}
    \leanok
    \uses{thm:detectability_lemma}
    In the setting of the previous theorem, let $\|\psi\|=1$ and
    $\varepsilon_\varphi=\langle\varphi|H|\varphi\rangle/\|\varphi\|^2$.
    Then $\|\varphi\|^2\le1/(\varepsilon_\varphi/g^2+1)$, where a quotient
    by zero is read as zero.
\end{theorem}

\begin{proof}
    \leanok
    If $\varphi\ne0$ and $g>0$, then
    $\|\varphi\|^2(\varepsilon_\varphi/g^2+1)
        =\langle\varphi|H|\varphi\rangle/g^2+\|\varphi\|^2\le1$. If
    $\varphi=0$ there is nothing to prove, and if $g=0$ the right side is
    $1\ge\|\varphi\|^2$.
\end{proof}

\begin{theorem}[Commuting projections]
    \label{thm:commuting_complement_product}
    \lean{ContinuousLinearMap.commute_complementProd,
        ContinuousLinearMap.isStarProjection_complementProd,
        ContinuousLinearMap.norm_sq_sub_norm_sq_complementProd_le}
    \leanok
    \uses{thm:complement_product_removed_mass}
    An operator that commutes with $h_0,\dots,h_{k-1}$ commutes with
    $Q_{k-1}\cdots Q_0$. If the $h_j$ commute pairwise, then
    $Q_{n-1}\cdots Q_0$ is an orthogonal projection and
    $\|\psi\|^2-\|\psi_n\|^2\le\langle\psi|\sum_{k<n}h_k|\psi\rangle$.
\end{theorem}

\begin{proof}
    \leanok
    A product of commuting orthogonal projections is an orthogonal
    projection. For the inequality, $h_k\psi_k=Q_{k-1}\cdots Q_0h_k\psi$
    has norm at most $\|h_k\psi\|$, and
    $\|h_k\psi\|^2=\langle\psi|h_k|\psi\rangle$.
\end{proof}

\begin{theorem}[Product of two projections]
    \label{thm:two_projection_product}
    \lean{ContinuousLinearMap.IsStarProjection.one_sub_three_mul_le_norm_sq_apply_apply}
    \leanok
    \uses{thm:star_projection_identities}
    For orthogonal projections $P,P'$ and a unit vector $x$,
    $\|PP'x\|^2\ge1-3(\|x-Px\|^2+\|x-P'x\|^2)$.
\end{theorem}

\begin{proof}
    \leanok
    Put $u=x-Px$ and $w=x-P'x$. Then
    $\langle x,PP'x\rangle=\langle x-u,x-w\rangle$, whose real part is
    $1-\|u\|^2-\|w\|^2+\operatorname{Re}\langle u,w\rangle
        \ge1-\tfrac32(\|u\|^2+\|w\|^2)$. This real part is at most
    $\|PP'x\|$, and squaring gives the claim when
    $1-\tfrac32(\|u\|^2+\|w\|^2)\ge0$; otherwise the right side is negative.
\end{proof}

\begin{theorem}[Converse of the detectability lemma]
    \label{thm:detectability_lemma_converse}
    \lean{ContinuousLinearMap.one_sub_three_mul_le_norm_sq_complementProd_comp}
    \leanok
    \uses{thm:commuting_complement_product, thm:two_projection_product}
    Let $(h_k)$ and $(h'_k)$ be two families of pairwise commuting
    orthogonal projections, $Q_k=1-h_k$, $Q'_k=1-h'_k$,
    $P=Q_{n-1}\cdots Q_0$ and $P'=Q'_{n'-1}\cdots Q'_0$, and $H=\sum_{k<n}h_k+\sum_{k<n'}h'_k$. Then
    $\|PP'\psi\|^2\ge1-3\langle\psi|H|\psi\rangle$ for every unit vector
    $\psi$.
\end{theorem}

\begin{proof}
    \leanok
    By the previous two theorems with $\|\psi-P\psi\|^2=1-\|P\psi\|^2$.
\end{proof}

For a nearest-neighbour chain whose terms are projectors, the even and the
odd terms form two such families, and $PP'$ is the operator $\mathrm{DL}(H)$ of
\cite[Definition def:DLO]{Arad2016RigorousRG}. The converse above holds for
every unit vector, not only for eigenvectors of $H$, and its bound
$1-3\langle\psi|H|\psi\rangle$ implies the printed bound
$1-4\langle\psi|H|\psi\rangle$.

\begin{theorem}[Order and kernels under range projection]
    \label{thm:support_projection_order_kernel}
    \lean{Matrix.IsHermitian.supportProj_eq_cfc}
    \lean{Matrix.IsHermitian.isStarProjection_supportProj}
    \lean{Matrix.PosSemidef.le_supportProj_of_norm_le_one}
    \lean{Matrix.PosSemidef.sum_supportProj_mulVec_eq_zero_iff}
    \leanok
    For a finite Hermitian matrix $A$, its range projection is
    $P_A=\mathbf1_{\mathbb R\setminus\{0\}}(A)$ and is an orthogonal projection.
    If $A\ge0$ and $\|A\|\le1$, then $A\le P_A$.
    For every finite family $h_i\ge0$,
    $\ker\sum_i h_i=\ker\sum_i P_{h_i}$.
    This is the range-projection reduction used in
    \cite[Section 6.1]{Arad2016RigorousRG}.
\end{theorem}
\begin{proof}\leanok
    In an orthonormal eigenbasis, $P_A$ has diagonal entries zero or one.
    On the spectrum of a positive contraction, $x\le\mathbf1_{x\ne0}$.
    Functional calculus gives the order comparison. A positive sum annihilates
    a vector precisely when every summand does. Each positive matrix and its
    range projection have the same kernel, proving the last identity.
\end{proof}
